\documentclass[../../main2.tex]{subfiles}

\begin{document}

	% ======================================================================
	% 骨架文件：小节标题与追问链为终版；正文由后续会话填充。
	% 扩增模式标注：本章 = 有压扩增 + 身份权重（部分守恒）。
	% 原书材料接入点：
	%   - 原书 part3 ch3（政治与权力结构）→ 本章主体
	%   - 原书 part2 ch2（生存策略分化）→ §11.4：集权 vs 碎片改写为
	%     两种有压扩增组织形态（从「地理结果」改为「扩增模式」）
	% ======================================================================

	\chapter{通道三：政治（经济 + 身份权重）}
	\label{ch:politics}

	\begin{motivation}
		上一章遗留的问题：经济交换在模型里是对称的等价交换——但交换者有身份、历史、偏好。当交换双方的身份关系变得重要，交换网络会变成什么？\\
		\textbf{核心动机}：论证政治从经济中分离的时刻——交换中加入不可货币化的偏好；权力的来源是对职能的集体认同；解释政治守恒性为何弱于经济。\\
		\textbf{扩增模式}：有压 + 部分守恒——资源基础守恒，认同可以凭空增减。
	\end{motivation}

	\begin{logicmap}
	\begin{tikzpicture}[node distance=0.9cm, every node/.style={draw, rectangle, rounded corners, align=center}]
	\node (q) {交换双方的身份变重要，会怎样？};
	\node (i) [below=of q] {从等价到不等价（§11.1）};
	\node (r) [below=of i] {认同赋予权力（§11.2）};
	\node (c) [below=of r] {权力的模因化 = 制度化（§11.3）};
	\node (s) [below=of c] {生存策略的两条扩增路径（§11.4）};
	\node (n) [below=of s] {权力需要约束（抛给第\ref{ch:law}章）};
	\draw[-{Stealth}] (q) -- (i);
	\draw[-{Stealth}] (i) -- (r);
	\draw[-{Stealth}] (r) -- (c);
	\draw[-{Stealth}] (c) -- (s);
	\draw[-{Stealth}] (s) -- (n);
	\end{tikzpicture}
	\end{logicmap}

	% ==================================================================
	\section{从等价到不等价}
	\label{sec:pol-nonequivalence}

	\textcolor{gray}{\textit{【骨架 · 追问链（终版）】交换中加入不可货币化的偏好 → 人情、义务、权力。这还是交换吗？→ 是，但定价失效：同一条烟在商店和岳父家的「价格」不同。}}

	\textcolor{gray}{（正文待写：标准句式——「如果交换只认货币尺度，定价就是全部信息。但现在加入约束：交换携带身份。这一步改变了所有性质——这就是政治从经济中分离出来的时刻。」同构对照：货币 $\leftrightarrow$ 权力/人情；定价 $\leftrightarrow$ 权力赋予/撤销；市场 $\leftrightarrow$ 政治博弈场。）}

	\begin{readingnote}
		\textcolor{gray}{（待写：你有没有想过，为什么给邻居送钱是侮辱，帮邻居搬家是情分？）}
	\end{readingnote}

	\paragraph*{逻辑衔接}
	\textcolor{gray}{（待写）定价失效之后，新的度量出现了——别人为什么听你的？下一节「认同赋予权力」。}

	% ==================================================================
	\section{认同赋予权力}
	\label{sec:pol-recognition}

	\textcolor{gray}{\textit{【骨架 · 追问链（终版）】权力的来源是「对职能的集体认同」。这解释了为什么政治守恒性弱于经济：认同可以凭空增加或消失。认同怎么会增加？→ 通过模因传播（回到第\ref{ch:meme}章——嵌套结构在此显形）。}}

	\textcolor{gray}{（正文待写：权力 = 被认同的命令权；认同的增减不受守恒约束（一场演讲可以让权力翻倍或清零）；军队、税收、信息控制是守恒的那部分底座。）}

	\begin{questionbox}
		\textcolor{gray}{（待写：反驳位——「枪杆子里出政权，认同有用吗？」回答：枪需要人扣扳机，扣扳机需要认同——纯暴力政权是最短命的政权（雇佣军的悖论）。）}
	\end{questionbox}

	\paragraph*{逻辑衔接}
	\textcolor{gray}{（待写）权力来自认同，但认同住在人脑里——人会死。权力怎么活得比人长？下一节「权力的模因化 = 制度化」。}

	% ==================================================================
	\section{权力的模因化：制度化}
	\label{sec:pol-institutionalization}

	\textcolor{gray}{\textit{【骨架 · 追问链（终版）】沿用原书：个人魅力 → 职能程序的编码，等同于模因化。为什么编码？→ 因为人会死，模因不死——编码才能跨代传递。}}

	\textcolor{gray}{（正文待写：魅力的信息损失问题 → 科层制 = 权力的冷却与压缩；宪法、礼仪、职称都是编码层；\lowfreq 标注「制度化 = 模因化」是同构类比而非恒等。）}

	\begin{reader_anticipation}
		\item \textcolor{gray}{（待写）编码的失真与僵化？→ 原书「结构与效率」的僵化论证在此接入；第\ref{ch:law}章：法律作为最重的编码。}
	\end{reader_anticipation}

	\paragraph*{逻辑衔接}
	\textcolor{gray}{（待写）制度化解释了权力怎么延续，还没解释权力怎么组织。下一节把原书 part2 ch2（集权 vs 碎片）接入新位置。}

	% ==================================================================
	\section{生存策略的两条扩增路径}
	\label{sec:pol-two-paths}

	\textcolor{gray}{\textit{【骨架 · 追问链（终版）】原书 part2 ch2 在此接入：集权 vs 碎片，改写为「两种不同的有压扩增组织形态」。哪个更好？→ 都不是——它们对应不同的可预测窗口与崩溃模式（v2 注：$\kappa$ 已随状态空间层退役，可预测窗口的白话版见第\ref{ch:prediction}章 §\ref{sec:pred-intervention} 与第\ref{ch:forecast}章）。}}

	\textcolor{gray}{（正文待写：集权 = 降内部交易成本、高 $\kappa$、快扩散快崩溃；碎片 = 冗余多、慢、抗单点故障；治水-集权 / 贸易-分权重述（第\ref{ch:motivation}章 §\ref{sec:mot-result} 的治水-贸易约束在此变成组织形态）。）}

	\begin{readingnote}
		\textcolor{gray}{（待写）}
	\end{readingnote}

	\paragraph*{逻辑衔接}
	\textcolor{gray}{（待写）无论哪条路径，权力都会自我扩张——这是它的扩增本性。谁来拦住它？下一章的入口：需要一个「即使权力方想作恶也无法」的机制。}

	% ==================================================================
	\section{本章小结与衔接}
	\label{sec:pol-summary}

	\textcolor{gray}{（正文待写：收束——政治 = 经济 + 身份权重；嵌套关系更新：文化 $\supset$ 经济 $\supset$ 政治。）}

	\paragraph*{逻辑衔接}
	\textcolor{gray}{（待写）第\ref{ch:law}章「通道四：法律（博弈结果的编码与定量方法）」：当动态博弈的结果有争议、且必须给出一个数字时，法律登场。}

\end{document}
